%% ma: L10-B15-0034
%% lop: 10
%% chuong: 6
%% bai: 15
%% dang: 10.15.2
%% loai: TL
%% muc: VD
%% dap_an: "-"
%% nguon: "Ảnh giáo viên gửi 21/11/2026: Chuyên đề VI – Toán 10 – Chương VI (Hàm số – đồ thị và ứng dụng), tr. 199–203, Dạng 2 Câu 5"
%% kien_thuc: "Lớp 10 Bài 15 (tập xác định); phân tích đa thức thành nhân tử (lớp 8)"
%% trang_thai: chinh-thuc
%% ghi_chu: "Dạng do giáo viên cung cấp (21/11/2026). Có thể xếp VDC nếu học sinh không thấy được phép phân tích; chọn mức thấp hơn."
\begin{ex}
Cho hàm số $y=\sqrt{x^4+4x^3+(m+5)x^2+4x+4+m}$. Tìm tất cả các giá trị của $m$ để hàm số xác định trên $\mathbb R$.
\loigiai{Phân tích: $x^4+4x^3+(m+5)x^2+4x+4+m=\left(x^2+1\right)\left[(x+2)^2+m\right]$. Vì $x^2+1>0$ nên cần $(x+2)^2+m\ge0$ với mọi $x$. Giá trị nhỏ nhất của $(x+2)^2+m$ là $m$ (tại $x=-2$). Vậy $m\ge0$.}
\end{ex}
